% ============================================================================
% page_layout.tex - Page Layout and Typography
% ============================================================================
% Demonstrates geometry, fancyhdr, multicol, setspace, and microtype
% for professional document layout.
%
% Compilation:
%   pdflatex page_layout.tex
% ============================================================================

\documentclass[12pt,a4paper,twoside]{article}

% ============================================================================
% PACKAGE IMPORTS
% ============================================================================

\usepackage[utf8]{inputenc}
\usepackage[T1]{fontenc}
\usepackage{lmodern}
\usepackage[english]{babel}

% --- geometry: precise control over page dimensions and margins ---
% For a twoside document the "inner" margin (binding side) is wider.
\usepackage[
    top=2.5cm,
    bottom=3cm,
    inner=3cm,       % left on odd pages, right on even pages
    outer=2cm,       % right on odd pages, left on even pages
    headheight=15pt, % must be tall enough for the header font
    includeheadfoot  % account for header/footer in total height
]{geometry}

% --- fancyhdr: custom headers and footers ---
\usepackage{fancyhdr}

% --- multicol: multi-column layout ---
\usepackage{multicol}

% --- setspace: control line spacing ---
\usepackage{setspace}

% --- microtype: subliminal refinements for better typography ---
% Enables character protrusion and font expansion for tighter, more even
% text with fewer hyphenations and overfull boxes.
\usepackage{microtype}

% --- Additional packages used in this document ---
\usepackage{lipsum}      % Filler text (lorem ipsum)
\usepackage{xcolor}
\usepackage{hyperref}
\hypersetup{
    colorlinks=true,
    linkcolor=blue,
    pdftitle={Page Layout and Typography},
    pdfauthor={LaTeX Student}
}

% ============================================================================
% FANCYHDR: DEFINE PAGE STYLES
% ============================================================================

% Clear all header/footer fields before customizing
\pagestyle{fancy}
\fancyhf{}

% Twoside layout: headers mirror on even/odd pages
%   LE = Left on Even pages,  RO = Right on Odd pages  (outer positions)
%   RE = Right on Even pages, LO = Left on Odd pages   (inner positions)

\fancyhead[LE,RO]{\textbf{Page Layout \& Typography}}  % Document title on outer
\fancyhead[RE,LO]{\textit{\nouppercase{\leftmark}}}    % Section name on inner
\fancyfoot[LE,RO]{\thepage}                            % Page number on outer
\fancyfoot[C]{\footnotesize LaTeX Student}             % Author name at center
\renewcommand{\headrulewidth}{0.4pt}
\renewcommand{\footrulewidth}{0.4pt}

% plain style (used on \maketitle and \tableofcontents pages)
\fancypagestyle{plain}{
    \fancyhf{}
    \fancyfoot[C]{\thepage}
    \renewcommand{\headrulewidth}{0pt}
    \renewcommand{\footrulewidth}{0pt}
}

% ============================================================================
% DOCUMENT METADATA
% ============================================================================

\title{\textbf{Page Layout and Typography}\\[0.4em]
       \large Using \texttt{geometry}, \texttt{fancyhdr},
              \texttt{multicol}, \texttt{setspace}, and \texttt{microtype}}
\author{LaTeX Student}
\date{\today}

% ============================================================================
% DOCUMENT BODY
% ============================================================================

\begin{document}

\maketitle
\tableofcontents
\clearpage

% ============================================================================
% SECTION 1: THE geometry PACKAGE
% ============================================================================

\section{Page Dimensions with \texttt{geometry}}
\label{sec:geometry}

The \texttt{geometry} package replaces \LaTeX's manual dimension settings
(\verb|\oddsidemargin|, \verb|\textheight|, etc.) with a single, readable
declaration. The current document uses:

\begin{verbatim}
\usepackage[
    top=2.5cm, bottom=3cm,
    inner=3cm, outer=2cm,
    headheight=15pt,
    includeheadfoot
]{geometry}
\end{verbatim}

\subsection{Common Presets}

\begin{description}
    \item[\texttt{margin=1in}]       All four margins equal to one inch
    \item[\texttt{a4paper}]          Set paper size to ISO A4 (210\,mm \texttimes{} 297\,mm)
    \item[\texttt{letterpaper}]      US Letter size (8.5\,in \texttimes{} 11\,in)
    \item[\texttt{twoside}]          Enable mirrored left/right margins for binding
    \item[\texttt{landscape}]        Swap width and height for landscape orientation
    \item[\texttt{includeheadfoot}]  Count header and footer within the total page height
\end{description}

\subsection{Useful Geometry Commands}

\begin{description}
    \item[\texttt{\textbackslash newgeometry\{...\}}]
        Change geometry mid-document (useful for landscape pages)
    \item[\texttt{\textbackslash restoregeometry}]
        Restore the geometry declared in the preamble
\end{description}

\lipsum[1]

% ============================================================================
% SECTION 2: THE fancyhdr PACKAGE
% ============================================================================

\section{Custom Headers and Footers with \texttt{fancyhdr}}
\label{sec:fancyhdr}

The \texttt{fancyhdr} package defines six header and footer slots:
\texttt{L}, \texttt{C}, \texttt{R} for left, center, right, combined with
\texttt{O}/\texttt{E} for odd/even pages in twoside documents.

\subsection{Basic Setup}

\begin{verbatim}
\pagestyle{fancy}
\fancyhf{}                               % clear defaults

\fancyhead[LE,RO]{Document Title}        % outer header
\fancyhead[RE,LO]{\leftmark}            % inner header (section name)
\fancyfoot[LE,RO]{\thepage}             % outer footer (page number)
\fancyfoot[C]{Author Name}              % center footer

\renewcommand{\headrulewidth}{0.4pt}    % horizontal rule below header
\renewcommand{\footrulewidth}{0.4pt}    % horizontal rule above footer
\end{verbatim}

\subsection{Available Marks}

\begin{description}
    \item[\texttt{\textbackslash leftmark}]  Current chapter name (in \texttt{book}/\texttt{report})
                                              or section name (in \texttt{article})
    \item[\texttt{\textbackslash rightmark}] Current section name (in \texttt{book}/\texttt{report})
                                              or subsection name (in \texttt{article})
    \item[\texttt{\textbackslash thepage}]   Current page number as a string
    \item[\texttt{\textbackslash today}]     Current date
\end{description}

\subsection{Custom Page Style for a Chapter}

\begin{verbatim}
\fancypagestyle{chapterstyle}{
    \fancyhf{}
    \fancyhead[C]{\textsc{Chapter \thechapter}}
    \fancyfoot[C]{\thepage}
    \renewcommand{\headrulewidth}{0.6pt}
}

% Apply at the start of a chapter:
\chapter{Introduction}
\thispagestyle{chapterstyle}
\end{verbatim}

\lipsum[2]

% ============================================================================
% SECTION 3: THE multicol PACKAGE
% ============================================================================

\section{Multi-Column Layout with \texttt{multicol}}
\label{sec:multicol}

The \texttt{multicol} package balances text across multiple columns on the
same page, unlike \LaTeX's native \texttt{twocolumn} mode which divides the
entire document. Columns can start and stop anywhere in the text flow.

\subsection{Two-Column Example}

The \verb|multicols| environment (note the \emph{s}) takes the number of
columns as its mandatory argument:

\begin{verbatim}
\begin{multicols}{2}
  \section*{First Story}
  Text for the left column ...

  \columnbreak           % Force a column break here

  \section*{Second Story}
  Text for the right column ...
\end{multicols}
\end{verbatim}

\medskip
\noindent Live demonstration with two balanced columns and a separator rule:

\setlength{\columnseprule}{0.4pt}   % rule between columns
\setlength{\columnsep}{20pt}        % space between columns
\begin{multicols}{2}

\subsection*{Column One: Getting Started}
\lipsum[3]

\columnbreak

\subsection*{Column Two: Going Further}
\lipsum[4]
\end{multicols}
\setlength{\columnseprule}{0pt}     % remove rule for rest of document

\subsection{Three-Column Layout}

Three columns are common in newsletters and reference cards:

\begin{multicols}{3}
\lipsum[5]
\end{multicols}

% ============================================================================
% SECTION 4: THE setspace PACKAGE
% ============================================================================

\section{Line Spacing with \texttt{setspace}}
\label{sec:setspace}

Many institutions require 1.5\texttimes{} or double line spacing for
submission. The \texttt{setspace} package provides convenient commands:

\subsection{Document-Wide Spacing}

\begin{verbatim}
\usepackage{setspace}

% In preamble — applies to whole document:
\singlespacing        % 1× (default)
\onehalfspacing       % 1.5×
\doublespacing        % 2×
\setstretch{1.3}      % arbitrary multiplier
\end{verbatim}

\subsection{Local Spacing}

The \texttt{spacing} environment (lowercase) applies spacing locally:

\begin{verbatim}
\begin{spacing}{1.5}
This paragraph has 1.5× line spacing only.
\end{spacing}

\begin{singlespace}
Footnotes and bibliography often revert to single spacing.
\end{singlespace}
\end{verbatim}

\noindent Live demonstration:

\begin{spacing}{1.8}
\lipsum[6]
\end{spacing}

Back to normal (single) spacing after the environment closes.

\lipsum[7]

% ============================================================================
% SECTION 5: THE microtype PACKAGE
% ============================================================================

\section{Improved Typography with \texttt{microtype}}
\label{sec:microtype}

The \texttt{microtype} package exploits micro-typographic extensions in
\texttt{pdflatex} (and partially in \texttt{lualatex}) to produce
noticeably better output with minimal configuration.

\subsection{What microtype Does}

\begin{description}
    \item[Character protrusion]
        Punctuation and certain letters (commas, periods, hyphens) are allowed
        to extend slightly into the margin, creating a visually straighter edge
        without altering the nominal column width.

    \item[Font expansion]
        Letters are stretched or shrunk by a tiny amount (default $\pm 2\%$)
        so that lines fit more evenly, reducing hyphenation and overfull
        \verb|\hbox| warnings.

    \item[Tracking]
        Uniform letter-spacing adjustments for small-caps or special headings.

    \item[Kerning]
        Fine-grained control over space between specific letter pairs.
\end{description}

\subsection{Usage}

Simply loading the package in the preamble activates all features
automatically for \texttt{pdflatex}:

\begin{verbatim}
\usepackage{microtype}
\end{verbatim}

No further configuration is needed for most documents. For advanced tuning:

\begin{verbatim}
\usepackage[
    protrusion=true,
    expansion=true,
    tracking=true,
    kerning=true,
    spacing=true
]{microtype}
\end{verbatim}

\medskip
The following two paragraphs illustrate the difference.
The first uses default settings (\texttt{microtype} already active);
the second is a verbatim note since disabling microtype mid-document
requires \texttt{\textbackslash microtypesetup} which is load-time only.

\lipsum[8]

\noindent\colorbox{yellow!30}{\parbox{0.97\linewidth}{\small
\textbf{Note:} With \texttt{microtype} loaded you rarely see overfull
\verb|\hbox| warnings.  Compare the output by commenting out
\verb|\usepackage{microtype}| in the preamble.}}

% ============================================================================
% SECTION 6: COMBINING LAYOUT PACKAGES
% ============================================================================

\section{Putting It All Together}
\label{sec:combined}

A typical professional document combines all five packages:

\begin{verbatim}
% 1. Margins (geometry)
\usepackage[a4paper, top=2.5cm, bottom=3cm,
            inner=3cm, outer=2cm]{geometry}

% 2. Headers and footers (fancyhdr)
\usepackage{fancyhdr}
\pagestyle{fancy}
\fancyhf{}
\fancyhead[LE,RO]{My Document}
\fancyhead[RE,LO]{\leftmark}
\fancyfoot[LE,RO]{\thepage}

% 3. Line spacing (setspace)
\usepackage{setspace}
\onehalfspacing

% 4. Multi-column (multicol) — used locally as needed
\usepackage{multicol}

% 5. Improved typography (microtype)
\usepackage{microtype}
\end{verbatim}

This combination produces a well-proportioned document that meets most
institutional submission requirements, improves readability, and presents
a professional appearance.

\end{document}

% ============================================================================
% COMPILATION NOTES
% ============================================================================
% pdflatex page_layout.tex
%
% Note: fancyhdr's \leftmark is populated from \section on the PREVIOUS page,
% so you may see blank headers on the first page of each section — this is
% normal behaviour with the fancy page style. Use \thispagestyle{plain} on
% section-start pages if needed.
% ============================================================================
